\documentclass{article}

% ============================================================
% MA 213 In-Class Worksheet Template
% Student-facing worksheet for lecture activities.
% ============================================================

\usepackage[margin=1in]{geometry}
\usepackage{amsmath}
\usepackage{xcolor}
\usepackage{parskip}
\usepackage{array}
\usepackage{enumitem}
\usepackage{booktabs}

\newcommand{\coursenumber}{MA 213}
\newcommand{\weekplaceholder}{15}
\newcommand{\lectureplaceholder}{35}
\newcommand{\lecturetitle}{Diagnosing Overdispersion: Think/Pair/Share}

\newcommand{\nameblank}{\rule{1.75in}{0.4pt}}
\newcommand{\idblank}{\rule{1.55in}{0.4pt}}
\newcommand{\shortblank}{\rule{1in}{0.4pt}}
\newcommand{\longblank}{\rule{0.93\linewidth}{0.4pt}}

\begin{document}

\begin{center}
  {\LARGE \coursenumber{} Week \weekplaceholder{}: Lecture \lectureplaceholder{}}\\
  {\large \lecturetitle{}}
\end{center}

\begin{enumerate}[label=\arabic*.,leftmargin=*,itemsep=.7em]
  \item \textbf{Your name:} \nameblank \hfill \textbf{BU ID:} \idblank
\end{enumerate}

\textbf{Big idea:} You just saw how the Negative Binomial's dispersion parameter $r$ changes its shape: $E(X)=\mu$, $Var(X) = \mu + \mu^2/r$, small $r$ = long right tail, large $r$ = close to Poisson. We haven't yet told you $\hat r$ for the social media likes data --- the lecture gets it by numerically maximizing the likelihood. Beat the reveal: get your own quick hand estimate of $r$ first, \emph{without} numerical optimization.

\section*{The Setup}

Recall the likes-per-post data from earlier ($n=200$ posts):
\[
\bar{x} = 21.9 \qquad s^2 = 151.8.
\]

\section*{Step 1: Think (3 mins)}

\textbf{(a)} Compute the variance-to-mean ratio $s^2/\bar{x}$ (also called the \emph{Fano Factor}). What would it be if the data were perfectly modeled by a Poisson distribution? Does a Poisson model look appropriate here? Why or why not?
\vspace{0.8in}

\textbf{(b)} A quick \emph{method-of-moments} estimate of $r$ plugs the sample mean and variance directly into $Var(X) = \mu + \mu^2/r$ and solves for $r$:
\[
s^2 = \bar{x} + \frac{\bar{x}^2}{r} \quad\Longrightarrow\quad r = \frac{\bar{x}^2}{s^2 - \bar{x}}.
\]
Compute $\hat{r}$ for the likes data.
\vspace{0.8in}

\newpage
\section*{Step 2: Pair (3 mins)}

\begin{enumerate}[label=\arabic*.,leftmargin=*,itemsep=.7em, start=2]
  \item \textbf{Partner's name:} \nameblank \hfill \textbf{BU ID:} \idblank
\end{enumerate}

Compare your answers with your partner. Then discuss:

\begin{itemize}
  \item Based on your $\hat r$, do you expect the fitted Negative Binomial to look highly overdispersed (long right tail) or close to Poisson (roughly symmetric bell shape)?
  \item \textbf{Bonus:} this method-of-moments $\hat r$ is quick to compute by hand, but the lecture is about to show you an $\hat r$ from numerically maximizing the likelihood instead. Do you expect the two to agree exactly? Why might they differ, even on the same data?
\end{itemize}

\textbf{Our thoughts:}
\vfill

\end{document}
